\chapter{Родитель тахионного знака клеточной incidence-моды}
\label{ch:v10-hopf-reservoir-cell-incidence-tachyonic-sign}

\section{Знак, который не меняется с ориентацией}

Для минимальной ориентированной границы
\begin{equation}
 B=\binom{-1}{1},\qquad L_B=BB^T
 =\begin{pmatrix}1&-1\\-1&1\end{pmatrix},\qquad
 \Spec L_B=\{0,2\}.
\end{equation}
Обращение ориентации даёт $B\mapsto-B$, но
\begin{equation}
 (-B)(-B)^T=BB^T.
\end{equation}
Любая положительная реберная метрика также сохраняет
неотрицательность $BWB^T$. Поэтому знак нельзя получить ни выбором
ориентации, ни положительной перенормировкой ребра.

\section{Дополнение Шура и исходная неустойчивость}

Устойчивый двухмодовый родитель
\begin{equation}
 H_+=\begin{pmatrix}2&1\\1&2\end{pmatrix},\qquad
 \Spec H_+=\{1,3\}
\end{equation}
после исключения вспомогательной моды имеет дополнение Шура
\begin{equation}
 2-\frac{1^2}{2}=\frac32>0.
\end{equation}
Напротив,
\begin{equation}
 H_-=\begin{pmatrix}2&2\\2&1\end{pmatrix},\qquad
 \det H_-=-2,
 \qquad 2-\frac{2^2}{1}=-2
\end{equation}
создаёт отрицательную эффективную моду лишь потому, что полный
квадратичный родитель уже имеет отрицательное собственное значение.
Квартальный член может ограничить потенциал снизу, но не объясняет
происхождение исходной неустойчивости.

\section{Канонический фермионный минус}

Фермионный вакуумный функционал простейшего симметричного дублета имеет
форму
\begin{equation}
 \Gamma_F(x)=-2\sqrt{1+x^2},\qquad
 \Gamma_F''(0)=-2.
\end{equation}
Здесь отрицательный знак следует не из ручной замены $L_B\mapsto-L_B$, а
из фермионного детерминанта. Это первый некруговой условный источник
тахионной восприимчивости. Однако требуемый фермионный носитель, его
бифундаментальная связь с клеточной относительной модой и порог
доминирования над положительной incidence-жёсткостью пока не унаследованы.

Аудит двенадцати маршрутов по шести критериям даёт матрицу полного ранга,
но
\begin{equation}
 N_{\rm pass}=0/12,\qquad
 N_{\rm inherited+derived+negative}=0/12,\qquad
 N_{\rm max}=5/6.
\end{equation}

\begin{theorem}[Знаковая граница клеточного родителя]
Ориентация и положительная реберная метрика не могут превратить
$BB^T$ в тахионный оператор. Устойчивое бозонное расширение также не даёт
отрицательного дополнения Шура. Отрицательная мода требует либо уже
неустойчивого бозонного родителя, либо статистического минуса фермионного
детерминанта.
\end{theorem}

\begin{nogocriterion}[Перенос неустойчивости не является её происхождением]
Если полный квадратичный родитель уже имеет отрицательное собственное
значение, интегрирование вспомогательной моды лишь переносит эту
неустойчивость в incidence-сектор. Такой шаг не закрывает origin без
независимого вывода отрицательного направления и его нормировки.
\end{nogocriterion}

\begin{protocol}\begin{enumerate}
\item проверено знаковых механизмов: \textbf{$12$};
\item ранг критериев: \textbf{$6/6$};
\item полных проходов: \textbf{$0/12$};
\item устойчивый бозонный incidence-источник: \textbf{$0/1$};
\item условный фермионный знак: \textbf{$1/1$};
\item физическое происхождение: \textbf{$0/4$};
\item следующий гейт: \textbf{общий носитель фермионной тахионной восприимчивости}.
\end{enumerate}\end{protocol}

Машинный сертификат сохранён в
\path{s2t_v10_cell_birth_four_volume_induced_newton_breathing_anomaly_hopf_reservoir_m4_cell_incidence_tachyonic_sign_parent_origin_gate_results.json}.